\section{相关工作}
\label{sec:related}

本文的核心命题是：用计算机体系结构（computer architecture）的类比框架来构想未来智能计算系统的完整分层设计。
这一命题并非在真空中产生——到2025年，已有大量工作从各自角度触及了"LLM系统需要操作系统式的抽象"、"Agent需要记忆管理"、"工具调用需要标准协议"等关键洞察。
本节按功能主题梳理最相关的文献，逐步揭示每条线索的贡献与局限，最终定位本文的独特贡献。

\subsection{LLM-as-OS / Agent Kernel：将大语言模型视为操作系统内核}
\label{sec:related-llmos}

最早系统化地提出"LLM与操作系统类比"的研究，是Ge等人在2023年的工作"LLM as OS, Agents as Apps"\cite{ge2023llmos}。
这项工作的核心洞察可以用一张系统映射图来理解：LLM对应OS kernel（负责核心推理与调度），context window（上下文窗口）对应内存，外部存储对应文件系统，硬件工具对应外设，软件工具对应编程库，用户提示对应用户命令。
换言之，他们将LLM不是一个"聊天机器人"，而是一个"内核进程"——这一比喻后来被Karpathy在社交媒体上进一步传播\cite{karpathy2023llmos}。
在本文的类比体系中，Ge等人的工作直接对应"第2层Agent OS"的概念萌芽。

Mei等人的AIOS\cite{mei2024aios}将这一思路从类比推进到可运行的系统实现。
AIOS的核心设计是将Agent运行时（runtime）中原本与应用逻辑耦合的职责剥离出来，放入一个专门的"kernel"中，包括：
(1) 调度器——管理多个Agent的并发执行；
(2) 上下文管理——在多个Agent间分配有限的context window；
(3) 记忆管理——维护短期与长期记忆的读写；
(4) 存储管理——管理文件与持久化数据；
(5) 访问控制——控制Agent对工具和数据的权限。
实验报告显示，这种解耦设计带来了最高2.1倍的运行加速。
在本文的类比体系中，AIOS实现了一个早期的Agent OS原型，其kernel职责划分直接对应传统OS中的进程管理、内存管理和文件系统模块。

AgentOS\cite{agentos2025}进一步将LLM定义为受"structured operating system logic"约束的"reasoning kernel"——即LLM不是自由的推理引擎，而是必须在操作系统逻辑框架内运行的受控推理核心。
这一视角强调的是LLM的推理能力需要被结构化的系统逻辑所约束和编排，而不是让模型无约束地自由生成。

到2025--2026年，这条研究线从学术原型向实际桌面和移动场景推进。
微软的UFO$^2$\cite{ufo2_2025}将HostAgent（负责理解用户意图）、AppAgent（负责操作特定应用）、GUI--API混合动作层和虚拟桌面并行操作组织成一个完整的"Desktop AgentOS"。
其创新在于同时支持GUI截图理解与原生Windows API调用两种交互方式，使得Agent既能像人一样"看"屏幕，也能像程序一样"调用接口"。
Aura\cite{aura2025}则面向移动场景，采用Hub-and-Spoke（枢纽-辐射）拓扑：一个特权System Agent负责解析用户意图，多个沙箱化的App Agent各自负责一个领域的任务执行。
Aura的设计哲学是"安全优先"——App Agent之间的通信必须经过System Agent的中介，类似于传统OS中进程间通信必须经过内核。

\textbf{与本文的关系。}
上述工作各自定义了LLM在系统中的角色（OS kernel / reasoning kernel / Agent Kernel），但尚未形成统一的分层模型。
本文在第\ref{sec:icam}节提出的ICAM模型将这些不同的"kernel"概念统一到一个6层体系结构中，并明确LLM应被定位为第1层（概率计算核心），而Agent OS是第2层的独立抽象——两者不应混为一谈。

\subsection{记忆层级与虚拟上下文管理}
\label{sec:related-memory}

大语言模型的context window（上下文窗口）容量有限——例如GPT-4 Turbo为128K tokens，Claude 3为200K tokens，Gemini 1.5 Pro可达1M tokens\cite{gemini2024}。
然而，无论模型窗口多长，实际可用容量始终是有限的，这与传统计算机中"物理内存容量有限"的问题完全同构。
因此，一系列工作借鉴操作系统中的\textbf{内存管理}技术来扩展Agent的可用上下文。

MemGPT\cite{memgpt2023}（后演化为Letta平台）是这一方向的奠基性工作。
其核心思路不是简单地加一个RAG（检索增强生成）\cite{lewis2020rag}来查找外部知识，而是借鉴传统操作系统中的\textbf{虚拟内存管理}机制：
将有限的context window类比为"主存"（main memory），将外部数据库类比为"磁盘"（disk），
通过一个类OS的内存管理器自动地在context window和外部存储之间"换入换出"（page in/page out）信息片段，
使得Agent获得一种"看上去拥有无限上下文"的能力。
MemGPT的控制流还显式使用了interrupts（中断）机制——当需要从外部存储检索信息时，触发一个中断将控制权交给memory manager。
这一设计直接对应本文类比中的"虚拟内存"概念。

MemOS\cite{memos2025}将问题提升到更系统的层次。
它将三种不同粒度的记忆统一为可管理的系统资源：
(1) plaintext memory——以文本形式存储的外部知识；
(2) activation-based memory——模型内部的隐层表示（类似CPU寄存器）；
(3) parameter-level memory——模型参数中编码的知识（类似固件/ROM）。
MemOS用"MemCube"作为封装单元，将内容与溯源信息（provenance）、版本管理（versioning）打包在一起，并试图在retrieval（检索）与parameter learning（参数学习）之间建立迁移桥梁。
这一工作对应本文类比中"从KV cache到参数存储的多级存储层次"。

MemoryOS\cite{memoryos2025}直接将OS内存管理原则应用于Agent记忆系统。
HiAgent\cite{hiagent2024}（发表于ACL 2025）则提出了"层次化工作记忆管理"——以子目标（subgoal）作为记忆分块单元，将长时任务的记忆组织成层次结构，类似于操作系统中的多级页表（multi-level page table）。
A-MEM\cite{amem2025}提出了"能动性记忆"（agentic memory）的概念——记忆不是被动存储的记录，而是可以由Agent自主组织、重组和更新的动态资源。
LongMemEval\cite{longmemeval2024}进一步说明，长期记忆不是"存更多聊天记录"这么简单，而是涉及信息抽取、时序推理、知识更新与拒答能力的复杂系统工程。

\textbf{与本文的关系。}
记忆子系统的工作已经高度系统化，但其讨论通常独立于Agent运行时（runtime）和I/O层，缺少全栈视角。
更重要的是，现有工作尚未明确将KV cache（模型推理时的键值缓存）作为类比中"硬件缓存"的对应物——而本文认为KV cache正是连接"概率计算核心"与"语义记忆层级"的关键硬件抽象。
本文在第\ref{sec:kv}节中详细阐述了这一对应关系。

\subsection{Agent框架与运行时}
\label{sec:related-agent-framework}

如果说前两个主题分别定义了"内核"和"内存"，那么本节聚焦于"程序是如何在内核上运行的"——即Agent的执行框架与运行时环境。

\textbf{基础能力：推理与行动的交织。}
ReAct\cite{yao2023react}是Agent执行范式的基础性工作。
在此之前，LLM的推理（reasoning，如Chain-of-Thought）和行动（acting，如调用工具）是分离的两个步骤。
ReAct的核心贡献是证明：将推理和行动\textbf{交织}（interleave）在同一流程中——即"思考一步，行动一步，观察结果，再思考"——可以显著提高任务完成质量。
这类似于传统编程中"CPU执行一条指令，访问一次内存，检查一个条件"的fetch-decode-execute循环。

Toolformer\cite{schick2023toolformer}解决了另一个基础问题：让模型自己学会\textbf{何时}调用API。
通过在训练数据中自动插入API调用示例，Toolformer使模型在生成过程中自主决定是否需要调用某个工具——这类似于操作系统中的"系统调用"（system call）机制，由用户程序主动发起请求。

\textbf{编程Agent的兴起。}
2024--2026年最引人注目的Agent进展集中在软件工程领域。
OpenAI的Codex CLI\cite{openai_codex_overview_2026}于2025年4月发布，是一个运行在终端中的Agent编程助手。
它能够读取代码仓库、编辑文件、执行命令，并在沙箱（sandbox）环境中安全运行\cite{openai_codex_sandbox_2026}。
2026年初，Codex CLI引入了subagent（子Agent）功能\cite{openai_codex_subagents_2026}——主Agent可以生成专门的子Agent并行执行子任务，每个子Agent拥有独立的上下文窗口、模型配置和沙箱。
这直接对应了传统OS中"父进程创建子进程"的fork模式。
Codex还支持通过AGENTS.md文件定义自定义指令\cite{openai_codex_agents_md_2026}和通过Agent Skills定义可复用的能力模块\cite{openai_codex_skills_2026}。

Anthropic的Claude Code\cite{anthropic_claudecode_overview_2026}是另一个代表性的Agent运行时。
与Codex类似，它是一个终端原生的编程Agent，能够编辑文件、执行shell命令、连接外部服务。
Claude Code的架构设计被学术研究系统地分析过\cite{dive_claudecode_2026}——研究发现其系统架构中仅约1.6\%是AI决策逻辑，其余98.4\%是运行时脚手架（operational harness），包括安全控制、工具编排、上下文管理等。
Claude Code也支持subagent\cite{anthropic_claudecode_subagents_2026}和Skills\cite{anthropic_claudecode_skills_2026}扩展，并通过CLAUDE.md文件实现项目级记忆\cite{anthropic_claudecode_memory_2026}。

OpenHands\cite{openhands_paper_2025}（前身OpenDevin）是一个开源的通用Agent平台，发表于ICLR 2025。
它在SWE-bench Verified基准上达到约70--74\%的问题解决率\cite{swebench2023}，支持多种LLM后端。
SWE-agent\cite{sweagent2024}则专注于软件工程任务，用更简洁的Agent-Computer Interface（ACI）设计实现了相近的性能。

\textbf{与本文的关系。}
这些Agent框架在架构上已经浮现出操作系统的特征：Codex的subagent对应进程管理，Claude Code的CLAUDE.md对应配置/环境管理，OpenHands的ACI对应系统调用接口。
但它们目前仍是各自独立设计的"应用层"工具，缺乏统一的分层抽象。
本文的ICAM模型为这些框架提供了统一的理论框架，将它们的功能映射到明确的体系结构层次中。

\subsection{工具调用与I/O协议}
\label{sec:related-io}

Agent需要与外部世界交互——查询数据库、调用API、操作文件系统。
这一需求在传统计算机中对应I/O子系统（输入/输出设备与总线）。
近两年的关键进展是这种I/O抽象正在从各自为政的工具调用走向\textbf{标准化协议}。

HuggingGPT\cite{hugginggpt2023}是早期的代表性工作：将LLM放到"controller"（控制器）位置，由它解析用户任务、选择合适的专家模型、协调执行并汇总结果。
这类似于操作系统中的设备驱动管理层——由OS统一管理不同设备（模型）的调用。
Gorilla\cite{patil2024gorilla}聚焦于一个更具体的问题：当API文档频繁更新时，模型如何保持调用的准确性？这对应了设备驱动版本兼容性的问题。

\textbf{MCP：模型到系统的I/O总线。}
2024年11月，Anthropic推出了Model Context Protocol（MCP，模型上下文协议）\cite{mcp_intro_2026}，定义为连接AI应用与外部数据源、工具、工作流的开放标准。
MCP采用客户端-服务器架构：MCP Client运行在AI应用侧，MCP Server封装具体的工具或数据源。
通过统一的JSON-RPC协议，任何支持MCP的AI应用都可以连接任何提供MCP Server的工具——Anthropic官方将其比喻为"AI的USB-C"。
到2026年，MCP生态已快速扩展，Claude Code\cite{anthropic_claudecode_mcp_2026}和众多第三方工具均支持MCP连接。

\textbf{A2A：Agent到Agent的互联协议。}
2025年4月，Google在Cloud Next '25大会上推出Agent-to-Agent Protocol（A2A）\cite{google_a2a_2025}，明确定位为与MCP互补的开放协议。
A2A解决的不是"模型到工具"的连接问题，而是"Agent到Agent"的协作问题——支持能力发现（capability discovery）、任务生命周期管理（task lifecycle）、长时任务和多模态协同。
A2A已于2025年捐赠给Linux基金会，成为行业标准候选。
A2A采用对等的客户端-服务器模型，每个Agent既可以是发起请求的客户端，也可以是响应请求的服务器。

\textbf{与本文的关系。}
将两者合起来看，MCP扮演模型到系统的I/O总线角色（类似USB/PCIe），A2A扮演Agent到Agent的互联协议角色（类似TCP/IP网络协议栈）。
在本文的ICAM模型中，这分别对应第4层（I/O与工具层）和第5层（互联与协议层）。
现有工作尚未将这两个协议纳入统一的分层抽象中，也没有在体系结构层面讨论其与上下文管理、内核调度等层的交互。

\subsection{多Agent协作与分布式视角}
\label{sec:related-multiagent}

当多个Agent协同工作时，系统呈现出分布式计算的特征——这与从单机到分布式系统的演进完全平行。

AutoGen\cite{wu2023autogen}（微软研究院，2023年）是多Agent系统的奠基性框架。
其核心设计是将"多Agent会话"（multi-agent conversation）作为通用基础设施——不同角色的Agent通过自然语言对话来协调任务。
AutoGen已演进到v0.4版本，并被整合进Microsoft Agent Framework中，支持可扩展的生产级多Agent系统。

MetaGPT\cite{metagpt2023}（ICLR 2024 Oral）引入了一个关键的设计思想：将人类组织中的\textbf{标准操作流程}（SOP, Standard Operating Procedure）显式编码到多Agent协作中。
具体而言，MetaGPT模拟了一个虚拟软件公司：产品经理写需求，架构师设计系统，工程师写代码，QA工程师测试——每个角色是一个Agent，SOP定义了它们之间的工作流程和信息传递方式。
这类似于分布式系统中的工作流编排（workflow orchestration），将无序的Agent交互约束为有序的流水线（assembly line）。

Internet of Agents（IoA）\cite{ioa2024}（发表于ICLR 2025）从"互联网"而非单机系统的视角来想象Agent生态。
其核心贡献是提出了一种Agent集成协议（agent integration protocol），支持异构Agent之间的动态组队、协作和分布式执行。
IoA的设计灵感来自互联网——正如互联网连接了不同厂商、不同操作系统的计算机，IoA要连接不同框架、不同公司开发的Agent。
这也催生了IETF层面的标准化努力（IoA Protocol Internet Draft）。

\textbf{与本文的关系。}
多Agent系统已经在研究对象层面浮现出"分布式系统"的形态：AutoGen对应消息传递机制，MetaGPT对应工作流编排，IoA对应网络互联协议。
但大多数工作仍停留在"框架和协议"层，未与单体Agent的内存管理、内核调度、I/O、权限控制一起纳入统一的体系结构。
本文的ICAM模型将多Agent协作映射到第5层（互联与协议层）和第6层（应用层），并讨论了跨层交互的设计原则。

\subsection{认知架构与安全治理}
\label{sec:related-cognition-security}

前面各节从系统功能角度（内核、内存、运行时、I/O、多Agent）梳理了相关工作。
本节转向两个更高层的主题：Agent系统应遵循什么样的认知架构？如何治理其安全性？

\textbf{认知架构。}
CoALA（Cognitive Architectures for Language Agents）\cite{sumers2024coala}从认知科学角度提供了一个理论框架。
CoALA将语言Agent组织为三个核心组件：
(1) 模块化记忆（modular memory）——分为短期工作记忆和长期记忆（语义记忆、情景记忆、程序记忆）；
(2) 结构化动作空间（structured action space）——内动作（推理、记忆检索）与外动作（工具调用、环境交互）；
(3) 通用决策过程（generalized decision-making process）——规划、执行、观察的循环。
CoALA的价值在于提供了"为什么这个系统不是纯软件栈拼装，而是真正的认知架构"的理论支点。
L2MAC\cite{holt2024l2mac}（ICLR 2024）则明确将其框架称为"first practical LLM-based stored-program automatic computer (von Neumann architecture)"，包含指令寄存器、文件存储、控制单元和独立的LLM Agent模块。
Mi等人\cite{mi2025llmagentscomputersystems}从"computer systems perspective"组织LLM Agent综述，明确表示受von Neumann architecture启发，把Agent拆成感知、认知、记忆、工具、行动等模块，并指出context window类似主存、数据库类似磁盘，而Agent目前缺少的正是"cache module"。

\textbf{安全治理。}
ArbiterOS\cite{arbiteros2025}指出了一个根本性问题：我们正在用\textbf{确定性软件工程}（deterministic software engineering）的思维模型来指挥\textbf{本质上概率性的处理器}（inherently probabilistic processors）。
这一矛盾要求在系统架构层面引入治理层。
ArbiterOS由此提出"Agentic Computer"心智模型——将LLM重新定义为"Probabilistic CPU"（概率CPU），并设计governor（治理器）、formal instruction set（形式化指令集）和HAL（硬件抽象层）等治理层概念。
在本文的类比体系中，ArbiterOS的核心洞察——模型是概率性的、需要确定性治理层——正是本文提出"双平面架构"（概率计算平面 + 确定性控制平面）的直接灵感来源之一。

CaMeL\cite{debenedetti2025camel}（Google DeepMind，2025）从安全角度将操作系统能力安全原则应用于LLM Agent。
其核心设计理念是"Don't execute data"——严格分离可信指令与不可信数据，通过提取控制流和数据流并分别施加安全策略来防御提示注入攻击。
IronClaw\cite{ironclaw2025}（2026）明确提出要像操作系统一样保护AI Agent——将OS中的权限分离、沙箱隔离、审计日志等安全机制迁移到Agent系统中。

\textbf{与本文的关系。}
认知架构线提供了理论根基，安全治理线提供了工程约束。
但这两条线之间缺乏统一——认知架构通常不讨论安全性，安全治理通常不讨论认知结构。
本文的双平面架构试图统一两者：概率计算平面承载认知能力，确定性控制平面承载安全治理，两者通过明确定义的接口交互。

\subsection{现有工作的共同边界与未解决问题}
\label{sec:related-gap}

上述六个主题共同逼近了本文设想的"大模型原生智能计算机"，但存在三个尚未被统一解决的关键问题。

\textbf{问题一：LLM到底扮演什么角色？}
现有工作对LLM在系统中的定位存在根本分歧。
Ge等人与AIOS将LLM放到OS核心位置\cite{ge2023llmos,mei2024aios}——LLM就是kernel本身；
ArbiterOS将其下沉为被kernel治理的"Probabilistic CPU"\cite{arbiteros2025}——LLM是被管理的处理器；
AgentOS将其定义为受structured OS logic约束的reasoning kernel\cite{agentos2025}——LLM是受约束的推理核心。
这三种定位导致了完全不同的系统设计。
本文提出的ICAM模型给出统一答案：LLM是第1层（概率计算核心/CPU），Agent OS是第2层（操作系统内核），两者是不同的抽象层次，不应混为一谈。
LLM提供计算能力，Agent OS提供管理与调度——正如CPU与OS kernel在传统体系结构中是两个独立但紧密协作的层次。

\textbf{问题二：缺乏全栈分层模型。}
没有哪篇主流工作同时将概率计算核心、语义记忆层级、工具与外设总线、Agent OS/kernel、互联协议、治理安全层、编程模型/指令语义、评价指标全部放在同一套体系结构中。
L2MAC\cite{holt2024l2mac}聚焦stored-program计算机模型，Mi等人\cite{mi2025llmagentscomputersystems}提供了von Neumann启发的综述框架，但各自只覆盖了体系结构的一个侧面。
记忆管理的工作不讨论I/O协议，Agent框架不讨论KV cache，安全治理不讨论多Agent协作——每条线的贡献都是真实的，但边界也是清晰的。

\textbf{问题三：概率性与确定性的矛盾未被架构化。}
ArbiterOS\cite{arbiteros2025}正确地指出了概率性处理器与确定性治理之间的矛盾，但只是提出了概念层面的"governance-first"范式。
CaMeL\cite{debenedetti2025camel}和IronClaw\cite{ironclaw2025}提供了具体的安全机制，但未将其嵌入完整的分层架构中。
本文通过"双平面架构"（dual-plane architecture）将这一矛盾架构化：概率计算平面承载LLM的推理能力，确定性控制平面承载OS的管理与安全职责，两个平面通过明确的接口协议协作。

正是这三个未解决的问题，定义了本文的贡献空间：提供一个统一的、分层的、双平面的体系结构模型，将上述六条线索整合在同一张总图中。